% Figaro Paper N: Protocol Extension and Runtime Composition
% Computer Science and Cryptography register. The paper covers what stands
% above the kernel as a research object — schema design as a CS discipline,
% the coordinator pattern formally, and runtime composition.
% Audience: software-architecture researchers, protocol-extension theorists,
% the COALA-adjacent protocol-design community, computer-science readers
% approaching the protocol from a systems-architecture vantage.
% Preprint, May 2026

\documentclass[11pt,a4paper]{article}

\usepackage[utf8]{inputenc}
\usepackage[T1]{fontenc}
\usepackage{amsmath,amssymb,amsthm}
\usepackage{mathtools}
\usepackage{booktabs}
\usepackage{listings}
\usepackage{graphicx}
\usepackage{hyperref}
\usepackage{cleveref}
\usepackage{enumitem}
\usepackage{xcolor}
\usepackage[margin=1in]{geometry}
\usepackage{natbib}
\bibliographystyle{plainnat}

\lstset{
  basicstyle=\ttfamily\footnotesize,
  breaklines=true,
  columns=fullflexible,
  keepspaces=true,
  showstringspaces=false,
  frame=single,
  framerule=0.4pt,
  rulecolor=\color{gray!30},
  xleftmargin=0pt,
  xrightmargin=0pt,
}

\newtheorem{theorem}{Theorem}[section]
\newtheorem{proposition}[theorem]{Proposition}
\newtheorem{corollary}[theorem]{Corollary}
\theoremstyle{definition}
\newtheorem{definition}[theorem]{Definition}
\newtheorem{example}[theorem]{Example}
\theoremstyle{remark}
\newtheorem{remark}[theorem]{Remark}

\newcommand{\figaro}{\textsc{Figaro}}
\newcommand{\schemaid}{\mathrm{schemaId}}

\title{%
  Protocol Extension and Runtime Composition:\\
  Schema Design, the Coordinator Pattern, and the Software Architecture
  Above the Kernel%
}

\author{
  Alessandro Daliana
}

\date{May 2026}

\begin{document}

\maketitle

% ════════════════════════════════════════════════════════════════════
\begin{abstract}
The \figaro{} kernel is a settlement primitive: two operations
(\texttt{commit} and \texttt{resolveProcess}), three storage mappings,
no upgrade path, no escape hatches. Its mechanism-design derivation
(asymmetric bonding, buyer dominance with atomic resolution, and the
escape-hatch-weakness theorem) and its multi-method formal
verification (TLA\textsuperscript{+}, symbolic execution,
property-based fuzzing, and Certora specifications discharging the
kernel's invariants) are treated separately and stay strictly at the
kernel tier.

This paper covers what stands above the kernel as a research object.
The kernel by itself is not a useful protocol; it is a settlement
primitive that becomes a protocol when extensions are composed onto
it. Three discipline questions arise. First, when should a protocol
extension be written, and what does it mean for an extension to
preserve the kernel's equilibrium properties? We treat the
\emph{protocol extension doctrine} as the answer: the anchored
artifact pattern, append-only identity, first-write-wins binding,
the boundary between per-instance payloads and shared reference
semantics. Second, how should new schemas be designed and verified
at the protocol layer? We treat \emph{schema design as a computer-science
discipline}: the four-layer verification stack (Layer~A
specification, Layer~B prover mirror [pending], Layer~C Solidity
validator, runtime-tier UI binding), the 18-schema catalogue as a
coordinated reference set, and the \texttt{SchemaRegistrationHelper}
atomic-bind pattern that closes the front-running window between
schema registration and validator binding. Third, how does the
\emph{coordinator pattern} preserve the kernel's bonding equilibrium
when an external mechanism is composed onto the kernel? We give a
formal definition of composition semantics, state sufficient
conditions for equilibrium preservation, and treat the
\texttt{AttestationCoordinator} as the worked verification example
where the parametric Core-immutability rules in Certora discharge
the conditions. Finally, we describe the \emph{runtime composition}
above the protocol tier: the seven-layer composition pipeline
(protocol truth, semantic derivation, institution assembly,
mechanism packages, service bindings, view definitions, skin
bundles), the assembly registry pattern, the module system, the
designer canvas, and the
\texttt{(marketing)}/\texttt{(app)} route-architecture split.

The paper is in computer-science register: software architecture as
a research object, with the discipline derived from the kernel's
security requirements rather than from convenience.
\end{abstract}

\medskip
\noindent\textbf{Keywords:}
protocol extension, schema design, coordinator pattern, equilibrium-preserving composition,
runtime architecture, software composition, content-addressed identity

% ════════════════════════════════════════════════════════════════════
\section{Introduction}
\label{sec:intro}

The \figaro{} kernel is a settlement primitive whose two operations,
three storage mappings, and frozen invariants produce a Nash
equilibrium under asymmetric bonding rules and buyer dominance with
atomic resolution. As deployed, the kernel satisfies those
invariants under multiple formal-verification methods. The kernel
itself stops at its boundary because that is where the formal
results live.

The kernel by itself is not a useful protocol. A settlement primitive
needs a graph above it that participants can compose: schemas typing
the agreement content, mechanism contracts coordinating specific
patterns, runtime surfaces letting humans and agents interact with
the resulting institutional shapes. The graph is what makes the
kernel applied; the kernel is what makes the graph trustworthy. The
two are inseparable in any working deployment, and the discipline
that governs the second half — how protocol extensions are written
and how runtime compositions are organized — is the subject of the
present paper.

The argument is in three parts.

\paragraph{The protocol extension doctrine.}
Extensions to the kernel must preserve the kernel's equilibrium
properties or they break the protocol. The discipline is therefore
not a matter of taste: there is a decision rule for whether a new
domain feature belongs in the protocol or stays in app logic, and
the rule is checkable. We state the rule, develop the anchored
artifact pattern as the recurring structural shape, and identify
the design guardrails the rule produces.

\paragraph{Schema design as a computer-science discipline.}
A \figaro{} schema is not a JSON document on a filesystem. It is a
typed-data definition deployed across four coordinated layers
(specification, prover mirror, on-chain validator, runtime-tier
binding) under append-only identity discipline, with first-write-wins
binding to the validator contract and content-addressed reference
to the specification document. The discipline is a computer-science
research object, not just an implementation detail, and we treat it
as such. The 18-schema catalogue currently shipped in the protocol
is the worked-example reference set.

\paragraph{The coordinator pattern formally.}
The \emph{coordinator pattern} is the discipline under which an
external mechanism contract composes with the kernel without breaking
the bonding equilibrium: an extension contract satisfies specified
conditions on its read/write profile against kernel state. We state
the conditions formally with composition semantics and prove that
the conditions are sufficient. The \texttt{AttestationCoordinator}
contract is the worked verification example, with the parametric
Core-immutability rules in
\texttt{certora/AttestationCoordinator.spec} discharging the
sufficient conditions.

\paragraph{Runtime composition.}
Above the protocol tier sits the runtime tier, where institutional
assemblies, mechanism packages, service bindings, view
definitions, and skin bundles compose into rendered institutions
that humans and agents interact with. We describe the seven-layer
composition pipeline, the assembly registry pattern, the module
system, the designer canvas, and the architectural separation
between marketing surfaces and operational surfaces. The runtime
is software architecture, not user-experience design; we treat it
as a CS object.

\paragraph{Paper organization.}
\Cref{sec:primitive} summarizes the settlement primitive in the
register the present paper uses.
\Cref{sec:doctrine} develops the protocol extension doctrine.
\Cref{sec:schemas} treats schema design as a CS discipline.
\Cref{sec:coordinator} formalizes the coordinator pattern.
\Cref{sec:runtime} describes runtime composition.
\Cref{sec:refusals} catalogues what the discipline refuses.
\Cref{sec:conclusion} concludes.

% ════════════════════════════════════════════════════════════════════
\section{The Settlement Primitive}
\label{sec:primitive}

The bonded primitive's mechanism-design derivation and its formal
verification are out of scope for the present paper. We summarize the
kernel features the present paper relies on.

\paragraph{Kernel surface.}
Two state-changing operations: \texttt{commit} (dual-signed EIP-712
commitment, locks asymmetric bonds) and \texttt{resolveProcess}
(buyer-only, atomic settlement of all orders in a process). Three
storage mappings: \texttt{processes} (\texttt{ProcessState}),
\texttt{orderStatus} (\texttt{uint8}), \texttt{orderProcessId}
(\texttt{bytes32}).

\paragraph{Six invariants.}
(i) asymmetric bonding (buyer locks $2P$, seller locks $2G$);
(ii) progressive collateralization across $N$-party process chains;
(iii) buyer dominance (only the root buyer can resolve);
(iv) atomic resolution (all active orders settle simultaneously or
not at all);
(v) immutable evidence (commits and resolutions emit unmodifiable
events);
(vi) no escape hatches (no admin, no timeout, no governance, no
unilateral exit).

\paragraph{Three tiers.}
\emph{Kernel}: the irreducible settlement primitive. \emph{Protocol}:
the kernel together with extensions composed under the discipline of
this paper. \emph{Runtime}: the protocol together with semantic
layers, builder surfaces, and rendered institutional shapes that
humans and agents interact with. The present paper takes the
three-tier distinction as given and develops the protocol and
runtime tiers.

% ════════════════════════════════════════════════════════════════════
\section{The Protocol Extension Doctrine}
\label{sec:doctrine}

The kernel is intentionally narrow. Its narrowness is what produces
the bonding equilibrium; widening it would weaken the equilibrium. Extensions to the kernel therefore must add capability
without widening the kernel itself. The doctrine answers the question
of how this is done.

\subsection{What the kernel secures}

At the kernel layer the protocol secures five things and only these:
(i) process topology (which orders compose into which processes);
(ii) economic obligations between counterparties (the bond posture
at each order); (iii) role-bearing order nodes (who is the buyer and
who is the seller at each order); (iv) lifecycle and settlement
history (when each order was committed and resolved); (v) atomic
process resolution semantics (the all-or-nothing settlement rule).

The kernel does not encode domain-specific meaning. A bonded
commitment between a passenger and an airline, between a shipper and
a forwarder, or between an agent and a service provider all reduce
to the same kernel objects (an order with a seller, a buyer, a
payment, and an agreement-hash). Domain meaning is added by
extensions that attach typed information to the secured process
graph.

\subsection{The anchored artifact pattern}

The recurring structural shape of a \figaro{} extension is the
\emph{anchored artifact pattern}. An anchored artifact family
exists when:

\begin{enumerate}[label=(\arabic*)]
  \item multiple parties or tools must share a stable interpretation
    of some referenced artifact;
  \item that interpretation must remain auditable over time;
  \item the protocol needs a minimal on-chain reference point for
    that artifact family.
\end{enumerate}

The pattern then is: off-chain semantics (the document, the
methodology, the field catalogue, the legal interpretation) plus an
on-chain anchor for shared reference integrity. The anchor carries
only what the chain needs to know to identify the artifact, verify
that the off-chain document referenced is the one that was committed,
and check that the artifact was admitted for use at the time of
reference. The anchor does not carry the semantics themselves.

\begin{definition}[Anchor record]
\label{def:anchor}
A protocol-layer anchor for an artifact family is a tuple
$\langle \schemaid, \mathrm{version}, \mathrm{uriHash},
\mathrm{admitted} \rangle$ where $\schemaid$ is the anchor identity,
$\mathrm{version}$ is the version within the family,
$\mathrm{uriHash}$ is an immutable cryptographic hash of the
URI pointing at the off-chain content, and $\mathrm{admitted}$ is
the binding's monotone admission flag (false until first
registration, true thereafter; the discipline forbids
deactivation).\footnote{In the deployed implementation
(\texttt{src/SchemaRegistry.sol}), the admission flag is stored as
\texttt{mapping(bytes32 \(\Rightarrow\) bool) public registered};
\texttt{version} and \texttt{uriHash} live on the
\texttt{SchemaRegistered} event log rather than as queryable storage.
The conceptual anchor is reconstructed by reading the event log under
the append-only-identity discipline of \Cref{sec:append-only}.}
\end{definition}

The minimal-anchor surface is intentional. Larger anchors would
freeze interpretive commitments into brittle on-chain state; smaller
anchors would lose reference integrity. The four-field shape is the
narrowest surface that supports cross-party agreement on which
artifact is being referenced and which version of it.

\subsection{Per-instance payloads versus shared reference semantics}

The discipline turns on a distinction the doctrine makes load-bearing:
the difference between \emph{per-instance payloads} (the data
attached to a particular order or process instance) and \emph{shared
reference semantics} (definitions whose meaning must remain stable
across parties, tools, or time).

Per-instance payloads are operational data values: a specific
delivery manifest, sealed address data, notes for a particular
fulfilment event. These are typically private, mutable at the
business level, or specific to one workflow instance. They do not
deserve a protocol-level anchor. They live as instance data on the
order or process that carries them.

Shared reference semantics are definitions whose interpretation must
remain stable across counterparties, tools, or time: a disclosure
schema, a manifest schema family, a certification framework
reference, a quality-assurance reference standard. These may justify
a protocol-level anchor under the anchored artifact pattern.

\subsection{The decision rule}

\begin{proposition}[Decision rule for protocol extension]
\label{prop:decision-rule}
A new domain feature belongs in the protocol if and only if all
three of the following hold:
\begin{enumerate}[label=(\arabic*)]
  \item the feature attaches meaning to the secured process graph
    rather than to private instance data;
  \item the feature requires stable shared interpretation across
    counterparties or tools;
  \item the protocol needs to preserve that reference integrity over
    time.
\end{enumerate}
A feature satisfying all three conditions is a candidate for the
anchored artifact pattern. A feature failing any of the three
belongs in app logic, off-chain infrastructure, or per-instance
payload handling.
\end{proposition}

The proposition is operationally checkable: each condition has a
yes/no answer for a given proposed feature, and a feature that fails
the check is not a protocol-extension candidate. The doctrine
treats the rule as binding: the discipline is the willingness to
keep features out of the protocol when the rule fails them, not just
the willingness to add features when the rule passes them.

\subsection{Design guardrails}

Four guardrails operationalize the rule:

\begin{enumerate}
  \item Do not push app-specific workflow logic into the kernel.
    App logic that belongs to one institution does not belong to the
    protocol, regardless of how compelling the institution is.
  \item Do not make the protocol so abstract that every document
    becomes a first-class ontology object. Bounded generality is the
    target; universal-ontology generality is a failure mode.
  \item Do preserve reusable patterns when multiple domains clearly
    need the same coordination primitive. Generality earned by
    multiple domain instances is the kind of generality the protocol
    should accommodate.
  \item Keep the protocol legible: process graph first, extension
    semantics second, app UX last. The order of legibility reflects
    the order of trust dependency.
\end{enumerate}

The guardrails do not contradict the decision rule; they are the
operational restatement of it. A proposed extension that fails any
guardrail typically also fails the decision rule, and vice versa.

\subsection{Bounded generality}

The discipline's central tension is between under- and over-generality.
Under-generality produces one-off app-specific modules that cannot
become reusable protocol concepts; the protocol fragments into
incompatible bilateral arrangements. Over-generality produces a fake
universal ontology, abstract registries disconnected from concrete
coordination problems, and protocol bloat that mistakes possibility
for scope; the protocol becomes a CS-academic exercise rather than
working coordination infrastructure.

The right level of generality is generic enough to support reusable
extension patterns but concrete enough to stay grounded in process
coordination, obligations, and verifiable reference integrity. The
test for this balance is not theoretical: a proposed extension is
the right level of generality if at least two distinct domains have
plausibly demonstrated need for the same coordination primitive, and
the primitive can be specified narrowly enough to support both
without becoming a universal-ontology object.

% ════════════════════════════════════════════════════════════════════
\section{Schema Design as a CS Discipline}
\label{sec:schemas}

A \figaro{} schema is the concrete realization of an anchored
artifact family. Designing a schema is the operational form of
applying the extension doctrine. We treat the discipline in
computer-science register: typed-data design, content-addressed
extension, verification-stack architecture.

\subsection{The four-layer verification stack}

A protocol-tier schema is deployed across four coordinated layers
that together discharge the schema's correctness:

\paragraph{Layer A: client-side specification.}
A canonical JSON specification of the schema, parsed by a TypeScript
meta-validator (a closed subset of JSON Schema covering string with
specified format, integer, bigint as decimal string, boolean, enum,
array, and object types with strict closed-schema rejection of
unknown fields). The Layer~A specification is the canonical
declaration; everything below it is derived from it. Per-stage
overrides express the staged-attestation pattern (an
\texttt{applied}-stage attestation may carry different fields than
an \texttt{inspected\_intact}-stage attestation against the same
schema). The specification also carries a
\texttt{categories} array of short tags
(\texttt{handoff}, \texttt{commerce}, \texttt{emissions},
\texttt{lifecycle}, \texttt{jurisdiction}, \texttt{evidence-law},
etc.) that runtime UIs use for discoverability and grouping; the tags
are non-normative for validation but form the discoverability surface
for a schema's runtime presentation.

\paragraph{Layer B: prover mirror (pending).}
A Rust implementation of the Layer~A validator running in the SP1
zero-knowledge prover that produces verifiable attestation receipts
during batched off-chain execution. The prover guest program mirrors
the TypeScript validator byte-for-byte to enforce the schema during
batched attestation execution. Layer~B is currently deferred; the
TypeScript validator (Layer~A) plus the on-chain validator (Layer~C)
serve as the conformance specification for the eventual Rust port.

\paragraph{Layer C: on-chain validator.}
A Solidity contract implementing \texttt{ISchemaValidator} that
ABI-decodes the schema's content (no on-chain JSON parsing; the
on-chain layer reads the canonical ABI encoding produced by Layer~A's
encoder) and reverts with typed custom errors on any violation. The
validator is pure (no external state reads, no \texttt{block.*} or
\texttt{tx.*}, no external calls) and binds to one schemaId via a
compile-time literal (\texttt{bytes32 public constant override
schemaId = keccak256("...")}); using a constructor-set
\texttt{immutable} schemaId would forfeit the EVM-enforced
determinism guarantee and is forbidden by the validator-contract
pattern.

\paragraph{Runtime-tier UI binding.}
A frontend integration that consumes the Layer~A specification via
\texttt{useSchemaValidator(schemaId)}, renders schema-typed forms
through the existing module and view system, and delivers the
ABI-encoded content to the chain through the Layer~C validator. The
runtime tier is not strictly part of the verification stack, but it
is part of the schema's deployment surface; an unintegrated schema
exists on chain and not in any product.

\paragraph{Lockstep contract.}
A new schema is not done until all four layers ship in lockstep. If
Layer~A says one thing and Layer~C says another, attestations
silently divide into two interpretations and the schema's reference
integrity collapses. The lockstep contract is operational discipline,
not a formal soundness theorem: Layer~A and Layer~C agreement is
enforced by the test suite and the deploy-script structure rather
than proven by the verification stack. The lockstep contract is
enforced at the deploy-script level (the schema's registration script
registers the \texttt{schemaId} and binds the validator atomically)
and at the
test-suite level (Foundry tests in \texttt{test/schemaValidators/}
cover well-formed input plus every typed-error revert path).

\subsection{Append-only identity}
\label{sec:append-only}

Schema identity is append-only. The discipline produces three
operational rules:

\begin{enumerate}
  \item No in-place rewriting of schema meaning. A registered schema
    cannot be modified through any contract operation; the schema
    record at $\schemaid$ is immutable from the moment of
    registration.
  \item No silent mutation of the content reference behind an
    existing identity. The \texttt{uriHash} field of the anchor is
    immutable; off-chain documents may be mirrored for availability
    but the canonical reference is the hash, not any URI.
  \item New meaning requires a new version or a new schema identity.
    A schema family that needs a substantive change registers a new
    \texttt{schemaId} (e.g., \texttt{figaro-foo-v1} $\to$
    \texttt{figaro-foo-v2}); never mutates the v1 contract or its
    Layer~A spec.
\end{enumerate}

The append-only rule preserves historical interpretability:
attestations filed under $\schemaid_1$ remain readable against the
exact schema anchor they were filed under, regardless of subsequent
schema-family evolution.

\subsection{First-write-wins binding}

The validator-contract binding to a schemaId is permissionless and
first-write-wins. \texttt{AttestationCoordinator.setValidator(schemaId,
validator)} is callable by any address and binds the validator
permanently for that schemaId. Subsequent calls revert with
\texttt{ValidatorAlreadySet}. The first-write-wins discipline is
deliberate: an admin-mediated binding would re-introduce governance
authority over the schema-validation surface; a permissionless
binding without first-write-wins would let any address overwrite an
existing binding and break in-flight attestations. The combination
is the only design that satisfies both no-admin and
no-overwrite.

\begin{remark}[Why first-write-wins is structurally necessary]
Consider the alternatives. \emph{Admin-mediated binding}: an admin
can rotate validators, but admin authority over the validation
surface re-introduces the discretionary actor the no-escape-hatches
discipline rules out. \emph{Mutable bindings}: any address can
rewrite a binding, breaking existing attestations under that
$\schemaid$. \emph{Time-locked bindings}: a delay buffers the
overwrite, but any non-zero time creates a window in which the
binding is uncertain. The first-write-wins rule produces a binding
that is both permissionless and stable: anyone can set it once;
no one can change it after.
\end{remark}

\subsection{The atomic-bind pattern}

The first-write-wins rule produces a front-running window between
\texttt{SchemaRegistry.registerSchema} and
\texttt{AttestationCoordinator.setValidator}. An adversary observing
a pending registration can race to set a malicious validator under
the new schemaId before the legitimate validator binds. The
malicious validator self-attests as bound to the schemaId
(satisfying the binding-integrity check) and then accepts content
the legitimate validator would have rejected.

The mitigation is atomic registration: both writes execute in a
single transaction. The 17 reference \texttt{figaro-*} validators
are bound this way at genesis (the 18th catalogue schema,
\texttt{figaro-topology-v1}, is manifest-only and has no on-chain
validator), with
\texttt{script/Deploy.s.sol:\_deployAndRegisterValidators}
composing the calls in a single deployment broadcast session
controlled by one signer (no third-party interleaving against the
deployer's nonce stream). For third-party schemas registered
post-deploy, the \texttt{SchemaRegistrationHelper} contract supplies
genuine single-transaction atomicity as a stateless, no-admin,
no-fee public helper:

\begin{lstlisting}
function registerSchemaAndValidator(
  bytes32 schemaId,
  uint64 version,
  bytes32 uriHash,
  address validator
) external;
\end{lstlisting}

The helper composes the two underlying public calls
(\texttt{SchemaRegistry.registerSchema} and
\texttt{AttestationCoordinator.setValidator}) atomically. The
helper has no authority over either underlying contract; it has no
admin, no fee, and no privilege. It is a permissionless composer
that makes the two-step pattern safe.

\subsection{The 18-schema catalogue}

The protocol currently ships eighteen schemas, treated as a
coordinated reference set rather than as independent extensions.

\begin{enumerate}
  \item \texttt{figaro-topology-v1} (manifest-only clause; no runtime
    validator)
  \item \texttt{figaro-handoff-v1} (physical-exchange mode)
  \item \texttt{figaro-commerce-v1} (currency, payment, line items)
  \item \texttt{figaro-geo-v1} (origin / destination geohash)
  \item \texttt{figaro-fulfilment-v1} (fulfilment method)
  \item \texttt{figaro-ghg-protocol-v1} (GHG Protocol Corporate
    Standard; Category-2 disclosure)
  \item \texttt{figaro-ghg-iso-14064-v1} (ISO 14064 family;
    Category-2 disclosure)
  \item \texttt{figaro-ghg-pas-2050-v1} (PAS 2050 product carbon
    footprint; Category-2 disclosure)
  \item \texttt{figaro-ghg-en-16258-v1} (EN 16258 transport-emissions
    methodology; Category-2 disclosure)
  \item \texttt{figaro-ghg-custom-v1} (custom / non-standard GHG
    methodology; Category-2 disclosure)
  \item \texttt{figaro-ghg-measurement-v1} (runtime grams CO2e;
    Category-1)
  \item \texttt{figaro-delivery-lifecycle-v1} (5-stage progression
    + evidence URI)
  \item \texttt{figaro-proximity-policy-v1} (committed detection
    band; Category-2)
  \item \texttt{figaro-proximity-proof-v1} (per-handoff signed
    witness; Category-1)
  \item \texttt{figaro-merchant-process-v1} (merchant per-role event
    enum)
  \item \texttt{figaro-courier-process-v1} (courier per-role event
    enum)
  \item \texttt{figaro-jurisdiction-v1} (off-chain forum / law /
    language)
  \item \texttt{figaro-consent-v1} (cryptographic acceptance of an
    off-chain document)
\end{enumerate}

The catalogue exhibits several patterns worth flagging. The
\emph{sister-schema pattern} (the five GHG schemas, the proximity
policy/proof pair, the merchant/courier process schemas) supplies
multiple specialized variants of a shared coordination concept;
content shape is shared, schemaId is per-standard, and per-standard
extensions can be added to a single sister without affecting
siblings. The \emph{committed-vs-runtime split} (proximity
policy/proof, GHG protocol/measurement) separates the band the
parties commit at agreement signing (Category-2, byte-equality
enforced) from the runtime witness data filed during execution
(Category-1, fresh per attestation). The \emph{manifest-only
clause} pattern (\texttt{figaro-topology-v1}) represents a
shared-vocabulary anchor whose enforcement is purely off-chain;
parties commit to the topology at contract-signing time, the
indexer reconstructs the DAG from event logs, and no runtime
attestation fires. Each pattern is a discipline-level choice the
catalogue's authors made deliberately, and each is reusable by
future schemas in the family.

\paragraph{Class A and Class B records.}
A useful evidentiary taxonomy distinguishes records the kernel
produces by construction from records participants attest under
schema discipline. \emph{Class~A} records are kernel byproducts
(commitment digest, bond deposits, cumulative-value accumulator
state, order status transitions, resolution events): emitted by
the kernel, schema-typed at the kernel level, discretion-free, and
cryptographically authenticated. \emph{Class~B} records are
discretionary schema-validated attestations (delivery confirmations,
GHG disclosures, handoff attestations, custody-of-seal events,
lifecycle stage transitions): produced by participants through the
\texttt{AttestationCoordinator}, schema-typed at the \emph{protocol}
level, validator-gated, and substantively contestable on
truth-of-content grounds. The schema-design discipline of this paper
produces Class~B records: Layer~A specifies the content shape,
Layer~C enforces the shape at the on-chain validator, the runtime
tier renders and submits, and the resulting attestation enters the
evidentiary record as Class~B input.

\subsection{What schema design is not}

Three classes of would-be extension fall outside the discipline:

\emph{Per-instance payload schemas} that try to make every
order-specific data shape into a registered schema. Most
order-specific data does not need a registered schema; it needs an
order-specific encoder and an off-chain decoder. The discipline
applies the decision rule (\Cref{prop:decision-rule}) at the
point of asking ``does this need stable shared interpretation
across parties or tools?'' If the answer is no, the data is a
payload, not a schema candidate.

\emph{Schemas with mutable freeform text or large on-chain
payloads} (typically larger than ~1KB). Mutable freeform text
violates append-only identity. Large payloads violate the
minimal-anchor surface. Both move the schema in the wrong direction.

\emph{Schemas with admin, pause, or upgrade hooks.} The
no-escape-hatches discipline applies to schemas as it applies to
the kernel. A schema's validator is pure, immutable, and
permissionless once bound; an admin-controlled validator
re-introduces the discretionary actor.

% ════════════════════════════════════════════════════════════════════
\section{The Coordinator Pattern}
\label{sec:coordinator}

We turn now to the second discipline question: under what conditions
does an external mechanism contract compose with the kernel without
breaking the bonding equilibrium? The discipline that answers this
is the \emph{coordinator pattern}: a set of sufficient conditions on
an extension contract's read/write profile against kernel state
under which the bonding equilibrium is preserved by composition. We
state the conditions formally with composition semantics and prove
that they are sufficient.

\subsection{Composition semantics}

Let $\mathcal{K}$ denote the \figaro{} kernel as a state machine with
state space $\mathcal{S}_{\mathcal{K}}$ (the values of the three
storage mappings: \texttt{processes}, \texttt{orderStatus},
\texttt{orderProcessId}) and transition function $\delta_{\mathcal{K}}$
restricted to the two operations \texttt{commit} and
\texttt{resolveProcess}. Let $\mathcal{M}$ denote an external
mechanism contract with state space $\mathcal{S}_{\mathcal{M}}$ and
transition function $\delta_{\mathcal{M}}$.

\begin{definition}[Composition]
\label{def:composition}
We extend the state space to include the kernel contract's
ERC-20 token-balance position $\beta_{\mathcal{K}}$ and the
event log $E_{\mathcal{K}}$ alongside $\mathcal{S}_{\mathcal{K}}$
proper. The composition $\mathcal{K} \otimes \mathcal{M}$ is the
joint state machine over the extended state space in which
$\mathcal{M}$ may read from $\mathcal{S}_{\mathcal{K}}$,
$\beta_{\mathcal{K}}$, and $E_{\mathcal{K}}$ but may not write to
any of them, and $\mathcal{K}$ neither reads from nor writes to
$\mathcal{S}_{\mathcal{M}}$. The composition's transition function
admits interleaving via the EVM call graph (an
$\mathcal{M}$ transition's call frame may invoke
$\mathcal{K}$'s public methods, including state-changing methods
like \texttt{commit}, but only subject to condition (i) of
\Cref{prop:coordinator} below: any such invocation either fails
authorization or executes as the calling party would have
executed it directly through $\mathcal{K}$'s public surface,
without producing $\mathcal{M}$-mediated state mutation that the
calling party did not authorize).
\end{definition}

The asymmetric read/write structure is the load-bearing constraint.
$\mathcal{M}$ may consult kernel state to make its decisions, but
$\mathcal{M}$ cannot modify kernel state. $\mathcal{K}$ does not
need to know about $\mathcal{M}$ at all; the kernel is
composition-blind. The extended state space is required to make
the composition definition match the operational semantics of
EVM contracts: invariants like progressive collateralization
involve $\beta_{\mathcal{K}}$ (the kernel's escrow balance) as
well as $\mathcal{S}_{\mathcal{K}}$, and the immutable-evidence
invariant is over $E_{\mathcal{K}}$ rather than purely over
state mappings.

\subsection{The coordinator pattern, formally}

\begin{proposition}[Equilibrium-preserving composition]
\label{prop:coordinator}
The composition $\mathcal{K} \otimes \mathcal{M}$ preserves the
kernel's bonding equilibrium if $\mathcal{M}$ satisfies all four of
the following conditions:
\begin{enumerate}[label=(\roman*)]
  \item \emph{Read-only on kernel state.} For every operation
    $\mathcal{M}.f$ and every kernel-state predicate $P$ that is
    invariant under $\delta_{\mathcal{K}}$, $P$ is invariant under
    $\delta_{\mathcal{M}}$ as well: that is, no $\mathcal{M}$
    operation can mutate kernel state. Formally: for every
    $\mathcal{M}.f$, the call frame on which $\mathcal{M}.f$
    executes does not invoke any kernel state-changing operation
    (\texttt{commit} or \texttt{resolveProcess}) on $\mathcal{K}$
    in a way that produces unauthorized state mutation.
  \item \emph{No alternative settlement path.} $\mathcal{M}$ does
    not provide an operation that produces value flows equivalent
    to those of \texttt{resolveProcess} but bypasses
    \texttt{resolveProcess} or modifies its preconditions.
  \item \emph{No discretionary lock-bypass.} $\mathcal{M}$ does not
    provide an operation that releases bonded capital from
    $\mathcal{K}$'s escrow under conditions different from those
    \texttt{resolveProcess} enforces.
  \item \emph{Validator-gated content (where applicable).} If
    $\mathcal{M}$ accepts content typed by a registered schemaId,
    $\mathcal{M}$ routes the content through the registered
    \texttt{ISchemaValidator} before accepting it; content typed by
    a schemaId without a registered validator is rejected.
\end{enumerate}
\end{proposition}

The conditions are stated in the claims they constrain: condition
(i) preserves the kernel state predicates the verification stack
establishes (status monotonicity, transition correctness, cumulative
integrity, etc.); condition (ii) preserves buyer dominance; condition
(iii) preserves the no-escape-hatches discipline; condition (iv)
preserves the schema-validation surface. A mechanism that satisfies
all four can be composed onto the kernel with the kernel's
equilibrium guarantees still in force.

\begin{proof}[Sketch]
The kernel's bonding equilibrium rests on the six invariants of
\Cref{sec:primitive}. Each of the four conditions above maps to
preservation of a subset of those invariants under composition.

Condition (i) preserves invariants (i)--(v): if
$\mathcal{M}$ cannot mutate kernel state, the kernel's existing
invariants remain invariants under composition. Formally, if
$P : \mathcal{S}_{\mathcal{K}} \to \mathrm{Bool}$ is invariant under
$\delta_{\mathcal{K}}$ (formally verified at the kernel tier), and
$\delta_{\mathcal{M}}$ does not produce state changes in
$\mathcal{S}_{\mathcal{K}}$, then $P$ is invariant under
$\delta_{\mathcal{K} \otimes \mathcal{M}}$ since the latter is the
disjoint union of $\delta_{\mathcal{K}}$ and $\delta_{\mathcal{M}}$.

Condition (ii) preserves invariant (iii) (buyer dominance) and
invariant (iv) (atomic resolution) by ruling out alternative paths
that would bypass them. If $\mathcal{M}$ provided an operation
producing the value flows of \texttt{resolveProcess} without
satisfying the buyer-only and all-or-nothing preconditions, the
escape-hatch-weakness theorem applies and the equilibrium is broken.

Condition (iii) preserves invariant (vi) (no escape hatches) by
ruling out lock-bypass operations.

Condition (iv) preserves the schema-validation surface as a
component of the protocol-tier integrity. A mechanism that accepts
unvalidated content under a schemaId would create a content surface
that downstream consumers cannot trust; the validation gate is
what makes schema-typed content protocol-grade rather than
mechanism-local.

The conditions are jointly sufficient: a mechanism satisfying all
four does not interact with kernel state in any way that could
break any kernel invariant or any schema-validation surface
property. The conditions are also \emph{necessary in the worst
case} (relaxing any one produces a documented anti-pattern): a
mechanism that mutates kernel state directly (relax condition i)
breaks the verified kernel invariants; a mechanism with an
alternative settlement path (relax condition ii) implements the
timeout / arbitrator / governance-override escape-hatch the
escape-hatch-weakness theorem forbids; a mechanism with a
discretionary lock-bypass (relax condition iii) implements an
admin-controlled fund-withdrawal path the no-escape-hatches
discipline forbids; a mechanism that accepts unvalidated schema-typed
content (relax condition iv) creates a content surface that
downstream consumers cannot trust as protocol-grade. We do not claim the conditions are necessary in
every model; weaker conditions on a mechanism with sufficiently
restricted power (e.g., a pure view contract) might suffice for
narrower preservation claims. The four conditions are sufficient
for full equilibrium preservation and necessary in the
worst-case sense above.
\end{proof}

\subsection{The AttestationCoordinator as worked example}

The \texttt{AttestationCoordinator} contract is the canonical worked
example of the coordinator pattern. The contract accepts attestations
against bonded orders, gates them through registered
\texttt{ISchemaValidator} contracts, and emits typed events without
ever modifying kernel state. A second admission gate enforces a
Merkle-inclusion proof against the order's signed
\texttt{agreementHash}: only clauses committed at contract-signing
time can be attested under, so a clause that was not part of the
bilateral agreement cannot land an attestation against the order.
The contract's verification surface discharges
\Cref{prop:coordinator}'s conditions explicitly.

\paragraph{Conditions discharged.}
The Certora specification \texttt{certora/AttestationCoordinator.spec}
contains seven declared CVL rules (eight including the parametric
expansion of one rule reported in the cloud verification artifact).
Two of them are parametric over every
public function on \texttt{AttestationCoordinator} and are the
direct discharge of condition~(i):
\texttt{attestationCannotChangeOrderStatus} and
\texttt{attestationCannotChangeProcessState}. Each rule is
universally quantified over methods $f$ on the contract: for any
public function $f$ and any kernel-state value $\sigma$ before and
after $f$'s execution, $\sigma$ is unchanged. The Certora prover
discharges this universally for the AttestationCoordinator's public
surface.

Conditions (ii) and (iii) are discharged by the contract's
construction: \texttt{AttestationCoordinator} has no operation that
transfers tokens from the kernel's escrow, and no operation that
calls \texttt{resolveProcess} or any kernel state-changing function
on the kernel. The contract is read-only on the kernel.

Condition (iv) is discharged by the
\texttt{noValidatorBlocksBuyerAttestation} CVL rule, which proves
that any attestation under a schemaId with no registered validator
reverts.

\paragraph{The four conditions thus reduce to verification-stack
requirements} that a candidate coordinator can be checked against
mechanically. The coordinator pattern is not a guideline to be
followed at the developer's discretion; it is a contract-level
discipline whose conditions are dischargeable with the existing
verification toolchain.

\subsection{What the pattern is not}

The pattern is not a generic external-call disclaimer. A mechanism
that calls into the kernel's view functions, reads state, and
returns derived values is composition-safe; a mechanism that calls
into the kernel's state-changing functions on behalf of a third
party may or may not be composition-safe, depending on whether the
third-party authorization is properly delegated. The
\texttt{IRoleResolver} interface supplies the authorization-delegation
pattern for the latter case (a role-based attestation routed through
a mechanism that consults the role-resolver to authorize the
caller).

The pattern is also not a substitute for the schema-validation
discipline. A mechanism that accepts schema-typed content must
route it through the registered validator; a mechanism that produces
schema-typed events must produce content the registered validator
would have accepted. The two disciplines (coordinator pattern and
schema design) are independent and both must be satisfied for an
extension to be protocol-grade.

% ════════════════════════════════════════════════════════════════════
\section{Runtime Composition}
\label{sec:runtime}

We turn now to the runtime tier: the layer above the protocol where
institutional shapes are composed and rendered. The runtime tier is
not part of the protocol's verification surface (it does not affect
kernel state and does not change the equilibrium), but it is part of
the protocol's deployment surface (without a runtime, the protocol
has no users, only contracts). The discipline at the runtime tier is
software architecture: how to compose institutional shapes from
reusable parts without re-implementing the protocol per deployment.

\subsection{Seven-layer composition pipeline}

The runtime composes institutional surfaces through seven layers.

\begin{enumerate}
  \item \textbf{Protocol truth.} Contract reads, event subscriptions,
    and the cumulative state derived from them. This layer is
    authoritative; nothing above it can override it.
  \item \textbf{Semantic derivation.} The derivation of role
    contexts, capability bindings, mechanism trust boundaries, and
    guarantee/risk objects from protocol state. This layer is also
    authoritative: a runtime that renders a role the semantic layer
    has not derived is rendering a role that is not in the protocol.
  \item \textbf{Institution assembly.} The structural declaration
    of an institutional shape: which roles compose, which mechanisms
    are bound, which views are exposed. The assembly is the runtime's
    canonical document.
  \item \textbf{Mechanism packages.} Reusable units that own a
    mechanism's contract bindings, semantic adapters, capability
    mappings, default modules, and guarantee/risk copy. A mechanism
    package is the unit the runtime composes; an assembly references
    packages rather than individual modules.
  \item \textbf{Service bindings.} Off-chain or hybrid infrastructure
    bindings: identity resolution, catalogue metadata, discovery and
    search, messaging and handoff, evidence transport, geospatial
    filtering. The bindings are resolved through stable interfaces,
    not hardwired per archetype.
  \item \textbf{View definitions.} Per-surface UI composition: route
    or surface id, accepted context, visible slots, module ordering,
    role-specific visibility.
  \item \textbf{Skin / branding fields.} Presentation-only
    personalization: theme tokens, imagery, type and spacing
    preferences, non-semantic copy. In the current implementation
    these are lightweight fields attached to the assembly
    binding (accent colour, logo, hero, theme class) rather than
    a heavyweight separate pipeline stage; we list them as a
    distinct layer because they sit at a distinct trust boundary
    (presentation-only, never semantic), not because they
    require their own composition machinery. Skin / branding
    fields cannot alter capability validity, mechanism
    authority, risk boundaries, or settlement semantics.
\end{enumerate}

The pipeline is layered: each layer consumes the layers below it
and produces output the layers above it consume. The runtime renders
an institution by walking the pipeline from protocol truth at the
bottom to skin bundle at the top. The structural property the
discipline enforces is that institution-specific mutation lives in
layers 3 through 7; the bottom two layers are runtime authority and
not institution-overridable.

\subsection{Assembly registry pattern}

An institution assembly is a JSON document declaring the
institutional shape: roles, mechanism packages, view definitions,
module placements, narrative defaults. The assembly is the runtime's
structural primary; the rest of the pipeline is derived from it. The
assembly registry pattern handles four operations:

\begin{itemize}
  \item \emph{Discovery}: which assemblies exist for a given
    archetype, role, or geographic region.
  \item \emph{Parsing}: validating the assembly's JSON against the
    runtime's schema for institutional shape.
  \item \emph{Validation}: cross-checking that the mechanism packages
    referenced exist, that the view definitions reference valid
    modules, and that the role mappings are coherent.
  \item \emph{Publication}: making a new assembly available to the
    runtime, with the assembly's content addressed for
    interoperability.
\end{itemize}

The pattern is implemented in
\texttt{frontend/lib/shared/assemblyRegistry.ts} (discovery and
indexing), \texttt{frontend/lib/shared/assemblyParser.ts} (parsing
and validation), and
\texttt{frontend/lib/shared/assemblyPublication.ts} (publication to
the runtime).

The current reference set comprises six assemblies:
\texttt{local-commerce} (the canonical food-delivery archetype),
\texttt{direct-sale} (a buyer-to-merchant single-leg arrangement),
\texttt{figaro-equipment-rental}, \texttt{figaro-freelance},
\texttt{figaro-procurement}, and
\texttt{figaro-disclosure-review}. Each is a worked example of the
assembly format and an operational template that can be specialized
into a deployed institution.

\subsection{Module system}

Modules are the unit the assembly's view definitions place in the
rendered surface. A module owns a UI component, a set of typed
actions it can dispatch, and the semantic context it consumes. The
module registry (\texttt{frontend/components/modules/}) catalogues
the modules currently shipped: catalogue editing, cart composition,
order creation, attestation submission, delivery coordination, GHG
disclosure, FIG token operations, and others.

The module system has a structural property worth flagging: a module
is composition-safe if and only if its dispatched actions go through
the typed action model (rather than producing direct contract calls)
and its semantic context comes from the semantic derivation layer
(rather than from local state hardcoded into the module). The
discipline matches the coordinator pattern's separation of concerns
at the runtime tier: the protocol-tier coordinator pattern says
external mechanisms read kernel state and emit events without writing
to the kernel; the runtime-tier module discipline says modules read
semantic state and dispatch typed actions without writing to the
runtime's authoritative layers.

\subsection{The lens system}

A particular runtime-tier construct worth treating separately is the
\emph{lens} pattern: a typed context object that a module receives
and that exposes a constrained subset of the runtime's full state.
A buyer-side module receives the buyer's lens onto the process; a
seller-side module receives the seller's lens; an auditor's module
receives the auditor's lens. Each lens exposes only what its role
context licenses; the module cannot reach beyond the lens to access
other roles' state.

The lens pattern is the runtime's enforcement of the role-segregation
that the kernel enforces at the bonding-and-resolution surface. The
kernel's role-segregation is mathematical (only the buyer can call
\texttt{resolveProcess}); the runtime's role-segregation is
operational (only a buyer-context module can dispatch buyer-side
actions). The two reinforce each other: a runtime that exposed
seller-side state to a buyer-context module would let the buyer
submit attestations as the seller, which the kernel would reject at
signature verification but which the runtime should not have
attempted in the first place.

\subsection{The designer canvas}

A specialized runtime surface deserves separate treatment: the
assembly designer canvas (\texttt{ProcessGraphCanvas} +
\texttt{AgreementDrawer}). The canvas is a composition surface
where assembly authors fork an existing reference assembly or
start from blank, modify the bonded-process graph, bind schemas
to per-node clauses through a side drawer, and save drafts to
local storage.

The canvas is not a wizard, not a recommendation engine, not a
generative-AI interface; it is a structural editor for assemblies.
The designer pulls the runtime's primitives (commerce-v1 commitments,
handoff-v1 attestations, fulfilment-v1 attestations, etc.) onto the
canvas as draggable elements; the author composes them into the
shape the institutional context requires; the resulting assembly is
saved as a draft.

The discipline the canvas enforces is structural: an assembly that
violates the runtime's composition rules cannot be saved. A
fan-out node where one fan-out lacks a corresponding mechanism
package is flagged; an attestation node referencing an unbound
schemaId is flagged; structural composition failures (mismatched
mechanism packages, dangling sub-orders, topology shapes the
runtime does not admit) are surfaced before the draft can be
persisted. The canvas surfaces the
discipline rather than hiding it; an author who saves an assembly
through the canvas has by construction produced one that satisfies
the runtime's structural integrity.

\subsection{The (marketing) / (app) split}

A practical runtime-architecture choice worth naming: the frontend
separates routes into two top-level groups,
\texttt{frontend/app/(marketing)/} for publication-and-explanation
surfaces and \texttt{frontend/app/(app)/} for transactional
surfaces. The split has three structural properties.

\emph{No wallet provider in marketing routes.} A user who has never
connected a wallet must be able to read every marketing route. The
\texttt{(marketing)} group does not load the wallet provider; the
\texttt{(app)} group does.

\emph{Wallet-connect is signing prerequisite, not login.}
\texttt{(app)} routes mount the wallet provider but do not gate
read-only content behind connection state. Inline write affordances
use the canonical \texttt{WalletGate} wrapper to require connection
at the moment of signing rather than at route entry.

\emph{Two distinct headers.} The marketing header carries marketing
nav (Discover, About, Cryptoeconomics, etc.); the (app) header
carries protocol-surface nav (Connect Wallet, Inbox, Orders, etc.).
The two are not collapsed.

The split's structural property is that a reader of the corpus's
publication content (the eight Zargham-discipline papers, the
\texttt{cryptoeconomics} page, the \texttt{groups} page, etc.) is
not asked to connect a wallet to read; a participant's
transactional surface is wallet-gated by structural design rather
than by per-page logic. The architectural separation matches the
trust-tier separation of the underlying protocol: marketing is
untrusted publication, (app) is signing-required interaction, and
the route architecture reflects the trust layers rather than
collapsing them into one shell.

% ════════════════════════════════════════════════════════════════════
\section{What the Discipline Refuses}
\label{sec:refusals}

The discipline is more visible in what it refuses than in what it
admits. We catalogue the recurring failure modes the discipline
prevents.

\paragraph{Kernel changes.}
The kernel is frozen. Any proposed extension that requires a kernel
change is, by the discipline, not a candidate extension; it is a
proposal for a different protocol. The escape-hatch-weakness theorem
makes this not a stylistic preference but a mathematical necessity:
any unilateral exit path from the bonded state weakens the Nash
equilibrium. The protocol's
narrowness is its security property.

\paragraph{Admin paths.}
Any proposed extension that introduces admin authority over schema
binding, validator selection, mechanism activation, or kernel
state is refused. The first-write-wins binding pattern is the
specific implementation of the no-admin discipline at the
schema-validation surface; analogous patterns apply elsewhere
(e.g., the \texttt{OperatorRegistry}'s permissionless registration
without admin gating).

\paragraph{Force-majeure escape clauses.}
Any proposed extension that lets a participant avoid bond loss in
external events (weather, regulatory action, force majeure) is
refused. The structural answer to force-majeure-style concerns is
composition with parallel insurance processes: the participant
buys insurance from an external insurer naming themselves as
beneficiary, and the insurance settlement is independent of the
bonded commitment. Bonded settlement does not contemplate
external-event-conditional release; insurance does. The two compose;
neither replaces the other.

\paragraph{Multi-currency cross-leg coordination.}
Any proposed extension that lets a process chain denominate
different legs in different currencies is refused. Multi-currency
coordination requires an external rate of substitution (oracle,
DEX, pre-agreed FX rate) that re-introduces a discretionary actor
and breaks the same-unit comparability on which the 2$\times$
Nash-stability result rests. Multi-currency coordination at the
participant level is achieved through wallet-side swaps before
commitment; the kernel sees one currency per process.

\paragraph{Schemas with mutable freeform text.}
Any proposed schema with large-text or mutable-text fields is
refused. The minimal-anchor surface and append-only identity
together rule out anchoring substantial mutable content on chain.
Such content belongs off-chain, content-addressed by hash with the
hash registered as the schema's reference.

\paragraph{Modules with direct contract calls.}
Any runtime-tier module that bypasses the typed action model and
produces direct contract calls is refused. The action model's
discipline is what makes the runtime's compositional safety
analyzable; a module that escapes the discipline can produce
state changes the runtime's other layers do not know about, which
breaks the composition pipeline.

\paragraph{Skin bundles that alter semantics.}
Any skin bundle that affects capability validity, mechanism
authority, risk boundaries, or settlement semantics is refused. A
skin is presentation-only; an institution that wants different
semantics writes a different assembly, not a different skin.

\paragraph{Soulbound reputation as protocol primitive.}
Any proposed extension that introduces a non-transferable
reputation token bound to a wallet at the protocol layer is
refused. Reputation is a graph-tier signal, not a kernel-tier
construct; introducing soulbound reputation at the protocol layer
reifies what is appropriately a graph-level pattern (settlement
history is already on chain and readable as reputation by any party
who wants to read it that way) into a privileged-credential
construct that weakens the kernel's actor-neutrality property.

\paragraph{Fee discounts conditioned on protocol-tier signals.}
Any proposed extension that conditions kernel-tier or
protocol-tier behavior on a fee discount, a green-bond
adjustment, or any other rate-modifying signal is refused. The
2$\times$ bonding ratio is the minimum viable equilibrium under
the minimum-multiplier result; introducing a fee-discount
mechanism that modifies the effective ratio for some class of
participant breaks the equilibrium for that class. Whatever
incentive a fee discount would supply belongs at the assembly
or operator-registry tier as a parameter the parties choose,
not at the bonding rule.

% ════════════════════════════════════════════════════════════════════
\section{Conclusion}
\label{sec:conclusion}

Above the kernel sits a protocol layer (extensions to the kernel
under the discipline of \Cref{sec:doctrine,sec:schemas,sec:coordinator})
and a runtime layer (institutional shapes composed under the
discipline of \Cref{sec:runtime}). Both layers are research objects
in computer-science register, and both have a discipline derived
from the kernel's security requirements rather than from convenience
or developer ergonomics.

The protocol-extension doctrine produces a decision rule
(\Cref{prop:decision-rule}) for whether a proposed feature belongs
in the protocol; schema design as a CS discipline produces a
four-layer verification stack with append-only identity and
first-write-wins binding; the coordinator pattern produces formal
sufficient conditions (\Cref{prop:coordinator}) for
equilibrium-preserving composition. Each is a contract the
discipline holds extension authors to; together they make extension
to the kernel a rigorous engineering practice rather than an
opinion-driven one.

The runtime-tier composition produces a seven-layer pipeline (from
protocol truth to skin bundle), an assembly registry pattern, a
module system constrained by the typed action model, a lens system
enforcing role-context separation, a designer canvas surfacing
discipline at authoring time, and a route-architecture split between
marketing and transactional surfaces. The runtime tier is not
verified the way the protocol tier is, but its discipline is
operationally checkable, and the open question is whether formal
verification at the runtime tier is worth the work.

The contribution is the discipline itself, made explicit and
research-grade. The present paper supplies the connecting tissue:
how the kernel becomes a protocol becomes a runtime, with the
discipline at each layer specified well enough that an extension
author or runtime architect can apply it.

\section*{Acknowledgments}
The protocol-extension and runtime-composition work developed
incrementally alongside the kernel. The extension doctrine document
originated as an architectural note on GHG schema anchoring and
generalized; the runtime model emerged from operational requirements
of the local-commerce reference archetype and generalized. The COALA
collaborators have been a consistent interlocutor on the legal-tier
interactions of the protocol-extension and schema-validation surfaces.

% ════════════════════════════════════════════════════════════════════
\begin{thebibliography}{99}

\bibitem[ERC-1271(2018)]{eip1271}
Giordano, F., Lockyer, M., and Castonguay, P.
\newblock EIP-1271: Standard Signature Validation Method for
  Contracts.
\newblock Ethereum Improvement Proposal 1271, July 2018.
\newblock \url{https://eips.ethereum.org/EIPS/eip-1271}

\bibitem[EIP-712(2017)]{eip712}
Joyce, R. and others.
\newblock EIP-712: Ethereum Typed Structured Data Hashing and
  Signing.
\newblock Ethereum Improvement Proposal, September 2017.

\bibitem[Solidity Team(2024)]{solidity_doc_2024}
Solidity Team.
\newblock \emph{Solidity Documentation}, version 0.8.x series.
\newblock Ongoing as of 2024.

\bibitem[OpenZeppelin(2024)]{openzeppelin_contracts_2024}
OpenZeppelin.
\newblock \emph{OpenZeppelin Contracts: Reference Implementations
  for Solidity Patterns}.
\newblock OpenZeppelin Contracts Documentation, ongoing as of 2024.

\bibitem[Certora(2024)]{certora_2024}
Certora.
\newblock \emph{Certora Verification Language (CVL): Documentation
  and Reference}.
\newblock Certora Technical Documentation, ongoing as of 2024.

\bibitem[Foundry(2024)]{foundry_2024}
Foundry.
\newblock \emph{Foundry: Ethereum Smart Contract Testing Framework}.
\newblock Foundry Book, ongoing as of 2024.

\end{thebibliography}

\end{document}
